\section{Literature review}\label{sec:literature_review}

\subsection{Investor sentiment}


Empirical asset pricing studies debate whether investor sentiment introduces transient noise or predicts fundamentals. \citet{delong1990noise} develop a theoretical model to demonstrate that unpredictable changes in noise trader beliefs could move prices away from fundamental values. To study how investor sentiment affects the cross section of stock returns, \citet{baker2006} construct a composite index from market proxies \footnote{The index includes discounts on closed end funds, NYSE turnover, the number of IPOs, returns on the first trading day, the equity share in new issues, and the dividend premium.} and show that high sentiment predicts lower returns, specifically for stocks that are difficult to value and arbitrage. Subsequent research extends these indirect measures to international markets \citep{baker2012global,schmeling2009} \footnote{\citet{baker2012global} construct composite sentiment indices for six international markets and show that global and local sentiment components predict future stock returns. \citet{schmeling2009} uses consumer confidence across 18 countries to measure investor sentiment and predict future returns.} and other asset classes including government bonds, commodities, and index options \citep{baker2012bonds,gao2015,han2008}.\footnote{\citet{baker2012bonds} find that stock market sentiment predicts the excess returns of government bonds. \citet{gao2015} construct a composite index from market proxies using the partial least squares method to explain the returns and comovements of commodity futures. \citet{han2008} show that market pessimism steepens the volatility smile and increases negative skewness for index options.} These indirect measures often capture unrelated market noise, making their predictive power dependent on the choice and aggregation of proxies \citep{huang2015}.

Textual analysis measures sentiment from expressed opinions rather than market outcomes, although model and sample choices still affect the measurement. Early textual analysis research on investor sentiment aligns with the transient noise perspective. \citet{tetlock2007giving} analyzes a daily market column from the Wall Street Journal and shows that media pessimism predicts price declines followed by reversals, which indicates temporary price pressure instead of fundamental information.


By measuring investor sentiment from text on social media, subsequent studies challenge the transient noise view. \citet{chen2014} analyze Seeking Alpha articles and comments to show that these opinions predict future stock returns and earnings surprises. This evidence suggests that social media platforms can aggregate dispersed information about firm value rather than merely generating random chatter. \citet{bartov2018} find similar results using Twitter data, showing that individual tweets posted before earnings announcements predict forthcoming earnings and returns. Recent research evaluates the economic magnitude of textual signals through portfolio construction. \citet{ke2021predicting} and \citet{lopez2023chatgpt} extract sentiment from financial news to form portfolios and document return predictability. \citet{ke2021predicting} show that the incorporation delay is longer for small stocks. \citet{lopez2023chatgpt} interpret subsequent drift through information processing constraints and limits to arbitrage.

One possible source of these conflicting findings is that sentiment and attention capture distinct dimensions of investor behavior. Traditional market proxies often conflate these two dimensions as they are based on the same aggregate trading outcomes. For example, \citet{baker2006} use market turnover to construct a sentiment index, whereas \citet{barber2008} and \citet{da2011} interpret abnormal trading volume and extreme returns as indirect proxies for investor attention. This overlap implies that trading metrics cannot distinguish the direction of investor sentiment from the intensity of attention or other liquidity factors. Textual data separate these dimensions by using linguistic tone to measure opinions and message volume to measure attention \citep{antweiler2004}. \citet{cookson2024} address this issue by analyzing messages from StockTwits, Twitter, and Seeking Alpha. By separating textual sentiment from message volume, they demonstrate that sentiment contains information relevant to stock prices, whereas attention reflects retail buying pressure and price overreaction.

\subsection{Managerial sentiment}

Managers possess private information about firm prospects, and investors use managerial disclosures to assess firm fundamentals \citep{beyer2010financial,healy2001information}. However, managerial tone can reflect both information and managers' subjective beliefs about firm prospects. On the one hand, tone reflects cognitive biases. Using an aggregate manager sentiment index, \citet{jiang2019} show that optimism predicts lower future aggregate returns and earnings surprises. They also document links between manager sentiment, extrapolation of recent performance, and overinvestment. \citet{chen2018morning} observe that the tone of earnings conference calls becomes more negative later in the day and that investors temporarily overreact to this variation. On the other hand, managerial tone can convey information about future economic conditions. At the macroeconomic level, \citet{demers2024aggregate} demonstrate that positive changes in earnings announcement tone predict economic growth in the following quarter. This evidence suggests that corporate tone also contains information relevant to firm value and future stock returns.

Specific details help investors evaluate whether fundamental information supports managerial tone.  \citet{crawford1982} develop a model to show that conflicts of interest between a sender and a receiver make vague communication unreliable. Since managers have incentives to exaggerate firm prospects, their optimism lacks credibility. In this case, investors tend to place more weight on managerial tone when specific and verifiable information supports the response. Using management earnings forecasts, \citet{hutton2003} report that positive forecasts provide information only when managers include verifiable statements about future performance. \citet{baginski2016verification} analyze management press releases and reveal that quantitative forecast news strengthens the market response to linguistic tone, whereas imprecision weakens the response. 

Managerial communication may also affect how investors perceive firm risk. \citet{gennaioli2015} develop a model of delegated asset management in which trust in fund managers reduces perceived risk and increases investor willingness to hold risky assets. \citet{germann2025trust} test this mechanism in an experimental setting and conclude that greater trust in a fund manager increases confidence and allocations to risky assets. These findings motivate testing whether the return association of managerial replies varies with the sentiment expressed in the preceding inquiry. \citet{loewenstein2001risk} argue that emotional states affect risk judgments. Specifically, anxiety shifts attention toward threats and increases an individual's tendency to seek and rely on external advice \citep{eysenck2007anxiety,gino2012}. These findings suggest that investors evaluate the same managerial reply differently depending on their initial sentiment state.

\subsection{Investor-manager interaction}

Interactive settings may provide corporate information beyond standard disclosures. First, interactive environments reveal additional fundamental details as they allow investors to publicly request specific information \citep{hollander2010}. In traditional disclosures, managers determine the content based on their internal expectations of investor demand. \citet{matsumoto2011} examine over 10,000 earnings calls from 2003 to 2005 to separate prepared presentations from discussion sessions. Using intraday absolute returns, they demonstrate that discussions contain more information than prepared presentations as external questions prompt managers to reveal new details.

Second, the way managers respond to explicit questions provides an additional pricing signal to the market. Response behavior affects prices through direct market reactions and subsequent information production. Investors may directly interpret a refusal to answer as a lack of transparency. \citet{hollander2010} examine individual requests during earnings calls and document that managerial refusal to answer questions is associated with negative intraday stock reactions. Rather than directly refusing to answer, managers often use evasive language to withhold sensitive information \citep{gow2021}. \citet{lee2016} evaluates how managers rely on scripted responses during discussions, reporting that this behavior predicts negative stock returns and lower future operating performance. Similarly, \citet{xie2023} show that boilerplate replies on Chinese platforms occur more frequently during weak firm performance and correlate with reduced market liquidity. Evasive responses can also restrict information production by external intermediaries. \citet{boustanifar2025} report that evasive replies precede less original reporting from professional media. With less substantive content to analyze, reduced media coverage restricts investor engagement and is associated with weaker abnormal returns.


Recent studies examine interactive platforms to understand their market implications. \citet{lee2022shall} analyze approximately 2.5 million investor questions and corporate replies on Chinese IIPs from 2010 to 2017. Using manual annotations and language algorithms, they evaluate the nature of investor questions and measure the level of interaction through platform activity. They find that most questions reflect the difficulties investors face when integrating existing public information. Their results show that greater platform activity is associated with higher market liquidity, narrower trading spreads, and greater price informativeness, suggesting that interactions lower information processing costs. Related research extends these platform activity metrics to corporate outcomes, showing a positive relationship with corporate investment efficiency \citep{xu2024}. 

We construct the Effective Sentiment Index (ESI) from matched investor inquiries and managerial replies on IIPs. The index accounts for the investor's preceding sentiment and whether the managerial response provides substantive information. We then test whether this interaction measure predicts stock returns beyond investor tone and managerial tone considered separately.



